% !TEX spellcheck = en-GB

\documentclass[11pt,a4paper]{article}
%\usepackage[obeyspaces]{url} \urlstyle{sf}%

\usepackage{xspace}
\usepackage{url}
%\let\bfseriesbis=\bfseries
%\renewcommand{\bfseries}{\sffamily\bfseriesbis}

\usepackage[top=2cm, bottom=2.5cm, left=2.5cm, right=2.5cm]{geometry}

\newcommand{\etc}{\emph{etc.}\@\xspace}
\newcommand{\ie}{\emph{i.e.}\@\xspace}
\sloppy

\begin{document}
%
\pagestyle{empty}


\title{Open Automata meet Choreographies}



\date{}

\maketitle

\section{General Information }
\begin{description}
%\item[Discipline principale:] Informatique
\item[Advisors :] Rabéa Ameur-Boulifa (Télécom Paris -- Sophia), Cinzia Di Giusto (UCA -- Sophia)  and Ludovic Henrio (CNRS -- ENS Lyon )
\item[Contacts: ] \url{rabea.ameur-boulifa@telecom-paris.fr}\\  \url{cinzia.di-giusto@univ-cotedazur.fr}\\ \url{ludovic.henrio@ens-lyon.fr}
\item[Who:] Master 1 or 2 Computer Science students
\item[Where:] Eurecom/ Laboratoire I3S 
%\item[Duration:] 5/6 months
\end{description}

\section{Topic}

\paragraph{Context.}


Over the past years, we have developed a theoretical framework for open systems, introducing the notion of open automata  \cite{arxiv-weakbisim,henrio:Forte2016,hou:hal-02406098} — labelled transition systems (LTSs) extended with parameters and holes representing unspecified components. Each transition is enriched with guards and assignments that capture complex data and behavioural dependencies between the automaton and its open interfaces. These models preserve key compositional properties such as bisimulation when composed.

Composition in open automata consists of filling holes with other automata, provided their interfaces are compatible. Currently, compatibility relies mainly on the syntactic matching of action labels and arities. However, this approach fails to capture the structured interactions—the sequences of inputs and outputs—occurring between components.

To overcome this limitation, we aim to annotate open automata with typing  information, that is, a "local" descriptions of their interaction protocols. These local views, expressed through behavioural types, capture how multiple participants cooperate. By linking open automata to choreographies, we can reason compositionally about interaction correctness, ensuring that local behaviours collectively conform to a coherent global protocol.

Hence, the central question becomes:
\begin{quote}
Can we synthesise a choreography (global type) from a set of interacting open automata to ensure safe and deadlock-free composition?\end{quote}

\paragraph{Objectives.}

The main goal of this internship is to connect the compositional semantics of open automata with the good properties of choreographies, local behaviors. Concretely, the work will follow these steps:

\begin{itemize}
\item {\bf State of the art:} Study the semantics of open automata and the theory of behavioural types \cite{Gay_Vasconcelos_2025}
\item {\bf Choreographic Compatibility:}
Design a compatibility condition for composing open automata based on choreographic specifications. This involves linking the local transitions of automata with the corresponding projections of a global type.

\item {\bf Correctness and Safety Properties:}
Use the choreographic structure to prove deadlock freedom and communication safety for composed systems.

\item {\bf Bisimulation and Equivalence:}
Investigate how choreographic equivalence (global behaviour preservation) aligns with existing notions of bisimulation between open automata.

\item {\bf Illustrative Case Studies:}
Develop examples and tool-based experiments showing how choreographies improve composability and verification.

\end{itemize}

Depending on the candidate’s profile and interests, the project may be approached from a theoretical perspective—by developing new formal foundations and proofs—or from a practical perspective, by extending existing tools or designing new ones to support the proposed framework.


\begin{thebibliography}{9}

\bibitem{Hans:2016}
Hans H\"uttel, Ivan Lanese, Vasco T. Vasconcelos, Lu\`is Caires, Marco Carbone, Pierre-Malo Deni\'elou, Dimitris Mostrous, Luca Padovani, Ant\`onio Ravara, Emilio Tuosto, Hugo Torres Vieira, Gianluigi Zavattaro:
\newblock {Foundations of Session Types and Behavioural Contracts}.
\newblock ACM Comput. Surv. 49(1): 3:1-3:36 (2016)



\bibitem{Gay_Vasconcelos_2025}
Gay, Simon J. and Vasconcelos, Vasco T.
\newblock{Session Types}.
\newblock{Cambridge University Press, 2025}


\bibitem{henrio:Forte2016}
Ludovic Henrio, Eric Madelaine, Min Zhang:
\newblock {A Theory for the Composition of Concurrent Processes}.
\newblock In Albert, E., Lanese, I., eds.: {36th International Conference on
  Formal Techniques for Distributed Objects, Components, and Systems (FORTE)}.
  Volume LNCS-9688 of Formal Techniques for Distributed Objects, Components,
  and Systems., Heraklion, Greece (June 2016)  175--194. 
\newblock  \url{https://hal.inria.fr/hal-01299562}

\bibitem{Fantechi:MLTCDP2019}
Alessandro Fantechi, Elie Najm, Jean-Bernard Stefani:
\newblock {From Behavioural Contracts to Session Types}.
\newblock Models, Languages, and Tools for Concurrent and Distributed Programming 2019: 278-297



\bibitem{hou:hal-02406098}
Zechen Hou, Eric Madelaine:
\newblock {Symbolic Bisimulation for Open and Parameterized Systems}.
\newblock In: {PEPM 2020 - ACM SIGPLAN Workshop on Partial Evaluation and
  Program Manipulation}, New-Orleans, United States (January 2020). 
 \newblock \url{https://hal.inria.fr/hal-02406098}

\bibitem{arxiv-weakbisim}
Rabéa Ameur-Boulifa, Ludovic Henrio, Eric Madelaine:
\newblock {Compositional equivalences based on Open pNets}.
\newblock Journal of Logical and Algebraic Methods in Programming, In press. (hal-03894031).
% \newblock \url{https://arxiv.org/abs/2007.10770}
\end{thebibliography}

\end{document}
